\documentclass{article}
\usepackage{amsfonts,amsthm,amsmath}
\usepackage[mathscr]{euscript}
\usepackage{enumitem}

\setcounter{section}{0}
\renewcommand{\thesection}{\S\arabic{section}}
\renewcommand{\thesubsection}{\arabic{section}}

\newtheorem{thm}{Theorem}
\newenvironment{theorem}[1]{
	\renewcommand{\thethm}{#1}
	\begin{thm}
		}{\end{thm}}

\newtheorem{lma}{Lemma}
\newenvironment{lemma}[1]{
	\renewcommand{\thelma}{#1}
	\begin{lma}
		}{\end{lma}}

\theoremstyle{definition}
\newtheorem{ex}{Exercise}
\newenvironment{exercise}[1]{
	\renewcommand{\theex}{#1}
	\begin{ex}
		}{\end{ex}}

\title{Chapter 1}
\author{Connor Mooneyhan}
\date{April 9, 2024}

\begin{document}

\maketitle

\begin{ex}
	A nonempty family of sets $\mathscr{R} \in \mathscr{P}(X)$ is called a \textbf{ring} if it is closed under finite unions and differences (i.e., if $E_1, \dots, E_n \in \mathscr{R}$, then $\cup _{1}^{n} E_j \in \mathscr{R}$, and if $E, F \in \mathscr{R}$, then $E \setminus F \in \mathscr{R}$).
	A ring that is closed under countable unions is called a \textbf{$\sigma$-ring}.
	\begin{enumerate}[label=\textbf{\alph*.}]
		\item Rings (resp. $\sigma$-rings) are closed under finite (resp. countable) intersections.
		\item If $\mathscr{R}$ is a ring (resp. $\sigma$-ring), then $\mathscr{R}$ is an algebra (resp. $\sigma$-algebra) iff $X \in \mathscr{R}$.
		\item If $\mathscr{R}$ is a $\sigma$-ring, then $\left\lbrace E \subset X : E \in \mathscr{R}\ \mathrm{or}\ E^c \in \mathscr{R} \right\rbrace$ is a $\sigma$-algebra.
		\item If $\mathscr{R}$ is a $\sigma$-ring, then $\left\lbrace E \subset X : E \cap F \in \mathscr{R} \ \mathrm{for\ all} \ F \in \mathscr{R} \right\rbrace$ is a $\sigma$-algebra.
	\end{enumerate}
\end{ex}
\begin{proof}
	\ \begin{enumerate}[label=\textbf{\alph*.}]
		\item If $\mathscr{R}$ is a $\sigma$-ring and $\left\lbrace E_n \right\rbrace_{n \in \mathbb{N}} \subset \mathscr{R}$, then \begin{equation*}
			      \bigcap_{n \in \mathbb{N}} E_n = \left( \bigcup_{n \in \mathbb{N}} E_n \right) \setminus \left( \bigcup \left\lbrace E_a \setminus E_b : a, b \in \mathbb{N} \right\rbrace \right) \tag{1}
		      \end{equation*}
		      Considering the left and right components of the difference in the right hand side of (1), the left is clearly a countable union of members of $\mathscr{R}$, and the right can be seen to be a countable union by identifying each $E_a \setminus E_b$ with $\frac{a}{b}$.
		      Thus the right hand side of (1) is a difference of countable unions of members of $\mathscr{R}$, so it is in $\mathscr{R}$.

		      If $\mathscr{R}$ is a ring, the argument above applies by making $\left\lbrace E_n \right\rbrace$ a finite set.

		\item In either case, if $\mathscr{R}$ is a ring (resp. $\sigma$-ring) and an algebra (resp. $\sigma$-algebra), then given $E \in \mathscr{R}$, $X = E \cup E^c \in \mathscr{R}$.

		      Conversely, if $\mathscr{R}$ is a ring (resp. $\sigma$-ring) and $X \in \mathscr{R}$, then given $E \in \mathscr{R}$, $E^c = X \setminus E \in \mathscr{R}$, so $\mathscr{R}$ is closed under complements.
		      $\mathscr{R}$ is already closed under finite (resp. countable) unions, so $\mathscr{R}$ is an algebra.

		\item Let $A = \left\lbrace E \subset X : E \in \mathscr{R}\ \mathrm{or}\ E^c \in \mathscr{R} \right\rbrace$.
		      Given $E \in A$, by definition $E^c \in A$ since $(E^c)^c = E \in A$.

		      Let $\left\lbrace E_n \right\rbrace_{n \in \mathbb{N}} \subset A$, and express this set as the union of two sets, namely \begin{equation*}
			      \left\lbrace E_n \right\rbrace_{n \in \mathbb{N}} = \left\lbrace F_j \right\rbrace_{j \in J} \cup \left\lbrace G_k \right\rbrace_{k \in K}
		      \end{equation*}
		      where $\left\lbrace F_j \right\rbrace_{j \in J} \subset A$ is the collection of all sets in $A$ who are members of $\mathscr{R}$ and $\left\lbrace G_k \right\rbrace_{k \in K} \subset A$ is the collection of sets in $A$ who are not themselves in $\mathscr{R}$, but whose complements are.
		      Then
		      \begin{align*}
			      \bigcup_{n \in \mathbb{N}} E_n & = \bigcup_{j \in J} F_j \cup \bigcup_{k \in K} G_k                            \\
			                                     & = \bigcup_{j \in J} F_j \cup \left( \bigcap_{k \in K} G_k^c \right)^c \tag{1}
		      \end{align*}
		      Since each $F_j$ is in $\mathscr{R}$, $\bigcup_{j \in J} F_j \in \mathscr{R} \subset A$.
		      Also, we know from (a) that $\sigma$-rings are closed under countable intersections, so since $G_k^c \in R$ for each $G_k$, we have $\bigcap_{k \in K}G_k^c \in R$.
		      Its complement is also in $R$ since $R$ is closed under complementation.
		      Thus (2) is a union of two members of $R$, and is therefore in $R \subset A$.

		\item Let $A = \left\lbrace E \subset X : E \cap F \in \mathscr{R} \ \mathrm{for\ all} \ F \in \mathscr{R} \right\rbrace$.
		      Note that $\mathscr{R} \subset A$ since $\mathscr{R}$ is closed under finite intersections.
		      Given any $E, F \in \mathscr{R}$, we have \begin{align*}
			      E^c \cap F = \left( X \setminus E \right) \cap F = F \setminus E
		      \end{align*}
		      which is a difference in $\mathscr{R}$ and is thus in $\mathscr{R} \subset A$.

		      Given $\left\lbrace E_j \right\rbrace_{j \in \mathbb{N}} \subset A$ and arbitrary $F \in A$, \begin{equation*}
			      \left( \bigcup_{j \in \mathbb{N}} E_j \right) \cap F = \bigcup_{j \in \mathbb{N}} \left( E_j \cap F \right),
		      \end{equation*}
		      which is a countable union of elements in $\mathscr{R}$.
		      Thus since the choice of $F$ was arbitrary, $\bigcup \left\lbrace E_j \right\rbrace \in A$.
	\end{enumerate}
\end{proof}

\begin{ex}
	Complete the proof of Proposition 1.2.
\end{ex}
\begin{proof}
	The construction of open intervals from half-open intervals and rays is standard.
\end{proof}

\begin{ex}
	Let $\mathscr{M}$ be an infinite $\sigma$-algebra.
	\begin{enumerate}[label=\textbf{\alph*.}]
		\item $\mathscr{M}$ contains an infinite sequence of disjoint nonempty sets.
		\item card($\mathscr{M}$) $\geq c$.
	\end{enumerate}
\end{ex}
\begin{proof}
	\ \begin{enumerate}[label=\textbf{\alph*.}]
		\item Choose a sequence $\left\lbrace E_n \right\rbrace_{n \in \mathbb{N}}$ of nonempty sets in $\mathscr{M}$ (this requires the axiom of choice).
		      We then derive the sequence $\left\lbrace F_n \right\rbrace$ defined recursively by $F_1 = E_1$ and $F_j = E_j \setminus F_{j - 1}$ when $j \geq 2$.
		      These are disjoint by construction and are all in $\mathscr{M}$ since $\mathscr{M}$ is closed under differences.
		\item Choose a sequence $\left\lbrace E_j \right\rbrace _{1}^{\infty}$ of disjoint nonempty sets in $\mathscr{M}$ (allowable by part (a) of this exercise).
		      Define $f: \mathscr{P}(\mathbb{N}) \to \mathscr{M}$ by $f: S \mapsto \bigcup_{k \in S} E_k$.
		      This is well-defined since $f(S)$ is a countable union of members of $\mathscr{M}$, so $f(S) \in \mathscr{M}$.
		      Also, $f$ is an injection since all the $E_j$ are disjoint.
		      This brings us to the conclusion that since $\mathrm{im }\,f \subset \mathscr{M}$, \begin{equation*}
			      \mathrm{card}\, \mathscr{M} \geq \mathrm{card}(\mathrm{im }\,f) = \mathrm{card}(\mathscr{P}(\mathbb{N})) = c.
		      \end{equation*}
	\end{enumerate}
\end{proof}

\begin{ex}
	An algebra $\mathscr{A}$ is a $\sigma$-algebra iff $\mathscr{A}$ is closed under countable increasing unions (i.e., if $\left\lbrace E_j \right\rbrace _{1}^{\infty} \subset \mathscr{A}$ and $E_1 \subset E_2 \subset \cdots,$ then $\bigcup _{1}^{\infty}E_j \in \mathscr{A}$).
\end{ex}
\begin{proof}
	If $\mathscr{A}$ is a $\sigma$-algebra, then $\mathscr{A}$ is closed under countable unions in general, so it is certainly closed under countable increasing unions in particular.
	On the other hand, if $\mathscr{A}$ is an algebra which is closed under countable increasing unions, then take any $\left\lbrace E_j \right\rbrace _{1}^{\infty} \subset \mathscr{A}$.
	Consider that \begin{equation*}
    \bigcup _{j = 1}^{\infty} E_j = \bigcup _{j = 1}^{\infty} \bigcup _{k = 1}^{j} E_k. \tag{1}
	\end{equation*}
  Clearly, for any $j \in \mathbb{N}$, $\left( \bigcup _{k = 1}^{j} E_k  \right)\subset \left( \bigcup _{k = 1}^{j + 1} E_k \right)$, and each $\bigcup_{k = 1}^{j}E_k$ is a finite union of elements of the algebra $\mathscr{A}$, and thus is in $\mathscr{A}$.
  Then the right hand side of (1) is a countable increasing union of members of $\mathscr{A}$.
  Thus (1) $\in \mathscr{A}$.
\end{proof}

\begin{ex}
  If $\mathscr{M}$ is the $\sigma$-algebra generated by $\mathscr{E}$, then $\mathscr{M}$ is the union of the $\sigma$-algebras generated by $\mathscr{F}$ as $\mathscr{F}$ ranges over all countable subsets of $\mathscr{E}$.
\end{ex}
\begin{proof}
  Let $F$ be the set of all countable subsets of $\mathscr{E}$.
  By Lemma 1.1, since $\mathscr{F} \subset \mathscr{E} \subset \mathscr{M}(\mathscr{E})$ for each $\mathscr{F} \in F$, $\mathscr{M}(\mathscr{F}) \subset \mathscr{M}(\mathscr{E})$.
  Their union, then, is also a subset of $\mathscr{E}$.

  We need to show, then, that $\mathscr{M}(\mathscr{E}) \subset \bigcup_{\mathscr{F} \in F}\mathscr{M}(\mathscr{F})$.
  If we can show that $\bigcup_{\mathscr{F} \in F} \mathscr{M}(\mathscr{F}) = \mathscr{M}(\bigcup_{\mathscr{F} \in F} \mathscr{F})$, then we will be done because $\mathscr{E} \subset \bigcup F$, so $\mathscr{M}(\mathscr{E}) \subset \mathscr{M}(\bigcup F)$.
  That is where we now turn our focus.

  Clearly $\mathscr{M}(\bigcup F) \subset \bigcup_{\mathscr{F} \in F} \mathscr{M}(\mathscr{F})$ since $\bigcup F \subset \bigcup_{\mathscr{F} \in F} \mathscr{M}(\mathscr{F})$, so we need to prove the other direction.
  Given $E \in \bigcup \mathscr{M}(F)$, fix some $\mathscr{F} \in F$ for which $E \in \mathscr{M}(\mathscr{F})$.
  Then since $\mathscr{F} \subset \bigcup F$, $\mathscr{M}(\mathscr{F}) \subset \mathscr{M}(\bigcup F)$, so $E \in \mathscr{M}(\bigcup F)$.
  As observed above, this completes our proof.
\end{proof}

\end{document}
